\section{Introduction}\label{sec:intro}

Regenerative chatter, whose linear mechanism is captured by stability lobe diagrams \citep{altintas1995analytical,insperger2002semi}, limits the material removal rate of most milling operations. For active damping, direct velocity feedback (DVF) through a proof-mass (inertial) actuator is attractive because the loop is collocated and stable up to a gain set by the actuator's own resonance \citep{elliott2001feedback}.

In a velocity feedback loop, a stroke-limited inertial actuator damps small vibrations and can lose stability at large ones \citep{baumann2007destabilization}. In milling the actuator must also react, through the feedback, to the forced, tooth-periodic vibration of the stable cut. \citet{ozsoy2025frequency} showed that actuator saturation can destabilise the control, and \citet{ozsoy2026robust} proposed a saturation-aware gain map: at each spindle speed, the DVF gain that maximises the linear stability limit subject to the peak force demand of a simulated cut staying within the actuator capacity. Both works concern the saturation caused by the cut that the design intends, not how much disturbance a cut that does not saturate can absorb before saturation sets in, and chatter follows.

That a linearly stable cut can chatter after a large enough perturbation is known \citep{kalmarnagy2001subcritical,stepan2005nonlinear}: loss of contact bounds the chatter amplitude \citep{tlusty1981basic}, linearly stable cuts coexist with chatter in bistable or unsafe zones \citep{dombovari2008estimates,dombovari2015bistable,dombovari2019experimental} whose basins of attraction have been computed \citep{yan2017basins}, and, with active damping, bistability has been studied in turning on a hardware-in-the-loop rig \citep{sahu2021bistability,sahu2023adaptive}. Invariant sets of predictive \citep{zhang2015ellipsoid} and adaptive \citep{wu2016adaptive} chatter controllers give regional guarantees under input constraints, and sector conditions give ellipsoidal estimates of the region of attraction of saturated systems through linear matrix inequalities (LMIs) \citep{hu2001book,gomesdasilva2005antiwindup,tarbouriech2011book}, also for periodic and delayed systems \citep{zhou2011periodic,li2017hinfinity,sakthivel2022jdsmc}. Above the uncontrolled stability limit no bounded input can stabilise the cut globally \citep{hu2001book}, but the size of its basin has not been measured, nor bounded, for a fixed active-damping law whose actuator saturates in milling.

This paper asks: how much disturbance can a linearly stable, actively damped milling cut tolerate before actuator saturation triggers chatter, and how should the gain schedule account for it? It answers for DVF with the command clipped at the actuator capacity, in a model of the case study of \citet{ozsoy2026robust}. The answer turns on the residual authority, the voltage $\Vmax-\bmax$ that the forced vibration leaves free: it sets the size of the basin of the forced orbit near the stability limit, and it can be budgeted. The contributions are:
\begin{enumerate}
\item The saturation-induced finite-amplitude instability of the clipped DVF loop is measured and explained (Section~\ref{sec:bistability}); an entry-independent authority condition (Section~\ref{sec:feasibility}) provides the forced orbit and its command. In published saturation tests on the same rig (Section~\ref{sec:model-second}), losses of stability between the two linear limits occurred only with saturation, but the model does not predict them.
\item The basin is estimated from inside (Section~\ref{sec:certificates}): for the averaged delay equation a circle criterion, with its unclipped forced orbit checked on a dense depth grid, gives global attractivity of the forced orbit below a depth whose median over the speeds is 0.97 of the uncontrolled limit; above it, LMIs in which the generalised sector condition is applied around the time-varying forced command, an adaptation of the usual reduction of the saturation level by the input already in use \citep{hu2001book,tarbouriech2011book}, give ellipsoidal inner estimates of the basin, numerically verified for collocation models of the averaged and of the time-periodic delay equation.
\item Basin-aware gain scheduling (Section~\ref{sec:basin-aware}) gives, at each speed and among the gains evaluated, the largest grid depth up to which every grid depth is linearly stable, keeps a voltage margin and has a certified kick, on the averaged collocation model, of at least a required instantaneous disturbance: a basin-aware counterpart of the stability lobe diagram.
\item A case study of a published delayed difference term that spends no authority on the forced orbit \citep{pyragas1992continuous,li2022ijamt}, and a Monte Carlo study of the residual authority under errors in the model (Section~\ref{sec:delayed}), show that linear margin is not tolerance of disturbances and how fragile the residual authority is.
\end{enumerate}
The tools are established: time-domain simulation of the cut with a force-limited actuator \citep{mancisidor2015coupled,ozsoy2026robust}, describing functions \citep{gelb1968describing}, a frequency-domain model of the clip \citep{ozsoy2025frequency}, the circle criterion \citep{zames1966inputII}, sector conditions \citep{gomesdasilva2005antiwindup} and collocation \citep{breda2005pseudospectral,butcher2009chebyshev}. The study is computational: at the published gains the model classifies 17 of the 19 published cutting tests of \citet{ozsoy2026robust} as measured (Section~\ref{sec:model-case}), and Section~\ref{sec:discussion} outlines how the predictions could be tested.
% src: results/cuts_main.json (21 marks, 19 distinct cuts; pred_paper equals measured at 17/19), results/paper_vs_ours.json (p09: n 19, agree_paper_gains 17); paper/notes/novelty.md (claims the article must not make)
